\subsection{Marcinkiewicz 乘子定理}
\label{sec:ch8-marcinkiewicz-multiplier}

由推论~\ref{cor:ch3-bounded-variation-multiplier}, 实直线上的有界变差函数是 $L^p$ 乘子. Marcinkiewicz 乘子定理表明, 可以把该假设削弱为二进区间上的条件.

\begin{Theorem}
\label{thm:ch8-marcinkiewicz-multiplier}
设 $m$ 是有界函数, 且在 $\mathbb R$ 的每个二进区间上具有一致有界的变差. 则 $m$ 是 $L^p(\mathbb R)$, $1<p<\infty$, 上的乘子.
\end{Theorem}

\begin{Proof}
证明类似推论~\ref{cor:ch3-bounded-variation-multiplier}. 给定二进区间 $I_j$, 令 $T_j$ 为与乘子 $m\chi_{I_j}$ 相伴的算子. 只考虑 $I_j=(2^j,2^{j+1})$, 另一情形完全相同. 对 $\xi\in I_j$,
\[
 (m\chi_{I_j})(\xi)=m(2^j)+\int_{2^j}^{\xi}dm(t).
\]
不失一般性可设 $m$ 在 $\mathbb R\setminus\{0\}$ 上右连续. 此时
\[
 T_jf(x)=m(2^j)S_jf(x)+\int_{2^j}^{2^{j+1}}S_{t,2^{j+1}}f(x)\,dm(t),
\]
其中 $S_{t,2^{j+1}}$ 是与乘子 $\chi_{[t,2^{j+1}]}$ 相伴的算子. 由 Hilbert 变换的调制公式和定理~\ref{thm:ch7-cz-weighted-strong}, 若 $w\in A_2$, 则 $S_j$ 与 $S_{t,2^{j+1}}$ 在 $L^2(w)$ 上一致有界. 由 Minkowski 不等式和假设, $T_j$ 也在 $L^2(w)$ 上有界, 常数只依赖于 $m$ 的 $L^\infty$ 范数和 $m$ 在 $I_j$ 上的全变差. 若 $T$ 是与 $m$ 相伴的算子, 则定理~\ref{thm:ch8-littlewood-paley} 和定理~\ref{thm:ch8-weighted-sequence-square} 给出
\[
\begin{aligned}
 \|Tf\|_p
 &\le C\left\|\left(\sum_j|T_jS_jf|^2\right)^{1/2}\right\|_p\\
 &\le C\left\|\left(\sum_j|S_jf|^2\right)^{1/2}\right\|_p
 \le C\|f\|_p.
\end{aligned}
\]
\end{Proof}

\hypertarget{chapter-08-page-179}{}
定理~\ref{thm:ch8-marcinkiewicz-multiplier} 可推广到高维; 先给出 $\mathbb R^2$ 的形式.

\begin{Theorem}
\label{thm:ch8-marcinkiewicz-multiplier-r2}
设 $m$ 是平面上的有界函数, 在 $\mathbb R^2$ 的每个象限中分别对两个变量二次可微, 且
\[
 \sup_j\int_{I_j}\left|\frac{\partial m}{\partial t_1}(t_1,t_2)\right|dt_1<\infty,
 \qquad
 \sup_j\int_{I_j}\left|\frac{\partial m}{\partial t_2}(t_1,t_2)\right|dt_2<\infty,
\]
以及
\[
 \sup_{i,j}\int_{I_i\times I_j}
 \left|\frac{\partial^2m}{\partial t_1\partial t_2}(t_1,t_2)\right|dt_1dt_2<\infty,
\]
其中 $I_i,I_j$ 是 $\mathbb R$ 中的二进区间. 则 $m$ 是 $L^p(\mathbb R^2)$, $1<p<\infty$, 上的乘子.
\end{Theorem}

\begin{Proof}
仍只考虑 $\mathbb R_+$ 中的二进区间. 令 $I_i=(2^i,2^{i+1})$, $I_j=(2^j,2^{j+1})$, 并固定 $(\xi_1,\xi_2)\in I_i\times I_j$. 则
\[
\begin{aligned}
 m(\xi_1,\xi_2)
 &=\int_{2^i}^{\xi_1}\int_{2^j}^{\xi_2}
 \frac{\partial^2m}{\partial t_1\partial t_2}(t_1,t_2)\,dt_1dt_2\\
 &\quad+\int_{2^i}^{\xi_1}\frac{\partial m}{\partial t_1}(t_1,2^j)\,dt_1
 +\int_{2^j}^{\xi_2}\frac{\partial m}{\partial t_2}(2^i,t_2)\,dt_2
 +m(2^i,2^j).
\end{aligned}
\]
若 $T_{i,j}$ 是与乘子 $m\chi_{I_i\times I_j}$ 相伴的算子, 则
\[
\begin{aligned}
 T_{i,j}f(x)
 &=\int_{I_i\times I_j}\frac{\partial^2m}{\partial t_1\partial t_2}(t_1,t_2)
 S_{t_1,2^{i+1}}^1S_{t_2,2^{j+1}}^2f(x)\,dt_1dt_2\\
 &\quad+\int_{I_i}\frac{\partial m}{\partial t_1}(t_1,2^j)S_{t_1,2^{i+1}}^1S_j^2f(x)\,dt_1\\
 &\quad+\int_{I_j}\frac{\partial m}{\partial t_2}(2^i,t_2)S_i^1S_{t_2,2^{j+1}}^2f(x)\,dt_2
 +m(2^i,2^j)S_i^1S_j^2f(x).
\end{aligned}
\]
上标表示算子作用的变量. 因此, 若 $w$ 在每个变量上一致满足 $A_2$ 条件, 即 $w\in A_2^*$, 则 $T_{i,j}$ 在 $L^2(w)$ 上有界, 见定理~\ref{thm:ch7-strong-maximal-weighted}. 于是可像定理~\ref{thm:ch8-marcinkiewicz-multiplier} 的证明那样得到结论.
\end{Proof}

\hypertarget{chapter-08-page-180}{}
定理~\ref{thm:ch8-marcinkiewicz-multiplier-r2} 的证明可立即推广到高维, 只是记号会很繁复. 在 $\mathbb R^n$ 中, 假设变为
\[
 \sup_{j_1,\ldots,j_k}
 \int_{I_{j_1}\times\cdots\times I_{j_k}}
 \left|\frac{\partial^km}{\partial\xi_{i_1}\cdots\partial\xi_{i_k}}(\xi)\right|
 d\xi_{i_1}\cdots d\xi_{i_k}<\infty,
\]
其中 $I_j$ 是实直线上的二进区间, 集合 $\{i_1,\ldots,i_k\}$ 遍历 $\{1,\ldots,n\}$ 的所有 $k$ 元子集, $1\le k\le n$.
